\chapter*{Preface}
\addcontentsline{toc}{chapter}{Preface}

This volume is an out-of-order history because experimental science was never
only a chronology. A date says when something happened. It does not say what
the record had learned how to do.

The route here follows a different order. It follows the obligations that an
experiment must satisfy before a claim can travel: a mark must be distinguished,
admitted, counted, encoded, compared, repeated, scaled, witnessed, localized,
compiled, and finally inferred. Those names come from the formal device used by
the earlier volumes. Volume 5 heard them as public language. This volume asks
where history learned them.

The first discipline of this book is realism. A famous experiment does not
enter because it is famous. It enters only if the apparatus changed what later
records were allowed to claim. A telescope matters only if the notebook, plate,
or witness protocol made private sight portable. A null result matters only if
the apparatus made absence into a record. A theory matters only after enough
records have compiled to license an inference that can be challenged without
being dissolved.

That is why the volume uses the structured \lean{phenom} form inherited from
the instrument book. Each episode must name a statement, an origin, an
observation, an operational constraint, and a consequence. The form is not
decoration. It is the guardrail against turning experimental history into
storytelling. If the observation cannot be named, the episode is not ready. If
the constraint cannot be stated, the experiment has not yet entered this book.
If the consequence says more than the record could carry at the time, the prose
has outrun the apparatus.

The second discipline is refusal. The best experiments in this book are not the
ones that confirmed what everyone hoped. They are the ones that made a prior
convenience unusable. The diagonal of a square refused the Pythagorean hope
that all magnitude could be carried by integer ratio. Galileo's telescope
refused naked-eye astronomy. Michelson and Morley refused ether drift as a
measured convenience. Planck's blackbody curve refused the old alphabet of
continuous energy exchange. N-rays and cold fusion refused prestige as a
substitute for repeatability. The trial registry refused the hidden negative
result. A discovery threshold refuses private excitement as public conclusion.

Refusal is not cynicism. A failed experiment can be bad work, but a disciplined
refusal is one of science's highest acts. It says: the world may be richer than
the record, but this record will not pretend to hold what it has not earned.

The third discipline is humility about truth. The experimental record can
license conclusions. It does not own the universe. This is not a softening of
science. It is the condition under which science remains public. A record that
knows its boundaries can be challenged, repeated, repaired, extended, and
trusted. A record that pretends to own the object has already stopped measuring.

So the book ends where the procession ends: \lean{INFERRED}. Not guessed. Not
felt. Not authorized by office or consensus. Inferred: licensed by a record
whose marks, refusals, transport conditions, and limits remain visible.

The order is not the calendar. The order is experimental necessity.
